\section{Kedudukan dan Peran Bahasa Assembly}
\label{sec:peran-assembly}

Prosesor pada akhirnya hanya mengenal sinyal listrik yang mewakili
nilai tinggi atau rendah. Seluruh instruksi yang dijalankan komputer
merupakan deretan bit yang disusun menurut aturan tertentu. Manusia
dapat menulis deretan bit itu secara langsung, tetapi cara tersebut
sangat lambat dan rawan kesalahan. Bahasa assembly hadir sebagai
representasi simbolik yang lebih mudah dibaca tanpa meninggalkan
kendali atas instruksi mesin.

Bagian ini diringkas dari \cite{sto2001}, Kata Pengantar, hlm.~3, dan
BAB I, hlm.~6--8, serta \cite{gunawan2019}, Bagian 1, hlm.~10--11.

\subsection{Tiga Tingkat Bahasa Pemrograman}

Berdasarkan kedekatannya dengan perangkat keras, bahasa pemrograman
dibagi menjadi tiga tingkat.

\begin{enumerate}
    \item \term{Bahasa mesin} adalah deretan bit yang langsung dikenali
          prosesor. Bahasa ini tidak memiliki nama simbolik, sehingga
          sulit dibaca dan ditulis manusia.

    \item \term{Bahasa assembly} adalah representasi simbolik dari
          bahasa mesin. Setiap instruksi mesin diberi nama pendek
          (\eng{mnemonic}) yang lebih mudah diingat, misalnya
          \lstasm|MOV| untuk memindahkan data. Hubungan antara
          assembly dan bahasa mesin bersifat \eng{one-to-one}: satu
          instruksi assembly diterjemahkan menjadi satu instruksi
          mesin.

    \item \term{Bahasa tingkat tinggi} seperti C, Java, atau Python
          menyembunyikan register dan memori di balik variabel,
          ekspresi, dan struktur kontrol. Satu baris bahasa tingkat
          tinggi dapat diterjemahkan menjadi banyak instruksi mesin.
\end{enumerate}

Hubungan ketiga tingkat tersebut ditunjukkan pada
Gambar~\ref{fig:tingkat-bahasa}.

\begin{figure}[H]
    \centering
    \begin{tikzpicture}[
        node distance=0.45cm,
        lapis/.style={
            draw=obeBlue,
            line width=0.8pt,
            rounded corners=3pt,
            minimum width=10.4cm,
            minimum height=1.05cm,
            align=left,
            inner sep=8pt,
            font=\small
        }
    ]
        \node[lapis, fill=obeBlueLight] (l4) {%
            \textbf{Bahasa tingkat tinggi}\\[2pt]
            \textit{satu baris} $\Rightarrow$ \textit{banyak instruksi mesin}};
        \node[lapis, fill=obeTealLight, below=of l4] (l3) {%
            \textbf{Bahasa assembly}\\[2pt]
            \textit{satu instruksi} $\Rightarrow$ \textit{satu instruksi mesin}};
        \node[lapis, fill=obeAmberLight, below=of l3] (l2) {%
            \textbf{Bahasa mesin}\\[2pt]
            \textit{deretan bit}};
        \node[lapis, fill=obeGrayLight, below=of l2] (l1) {%
            \textbf{Perangkat keras}\\[2pt]
            \textit{gerbang logika, bus, register}};

        \draw[-{Stealth[length=2.5mm]}, obeBlue, line width=0.8pt]
            ([xshift=-0.7cm]l4.north west) --
            ([xshift=-0.7cm]l1.south west)
            node[midway, left, font=\scriptsize, align=center]
            {semakin\\dekat\\mesin};
        \draw[-{Stealth[length=2.5mm]}, obeBlue, line width=0.8pt]
            ([xshift=0.7cm]l1.south east) --
            ([xshift=0.7cm]l4.north east)
            node[midway, right, font=\scriptsize, align=center]
            {semakin\\tinggi\\abstraksi};
    \end{tikzpicture}
    \caption{Tingkat abstraksi bahasa pemrograman, dari perangkat
    keras hingga bahasa tingkat tinggi.}
    \label{fig:tingkat-bahasa}
\end{figure}

Perbedaan ketiga tingkat tersebut dirangkum pada
Tabel~\ref{tab:perbandingan-bahasa}.

\begin{table}[H]
    \centering
    \caption{Perbandingan bahasa mesin, assembly, dan bahasa tingkat
    tinggi.}
    \label{tab:perbandingan-bahasa}
    \small
    \begin{tabular}{@{}p{2.9cm}p{3.8cm}p{3.8cm}p{3.6cm}@{}}
        \toprule
        \textbf{Aspek} & \textbf{Bahasa Mesin} &
        \textbf{Bahasa Assembly} & \textbf{Bahasa Tingkat Tinggi} \\
        \midrule
        Satuan penulisan & Deretan bit & Satu \eng{mnemonic} per
        instruksi & Ekspresi dan pernyataan \\
        \addlinespace
        Contoh & deretan bit & \lstasm|MOV AX, 5| & \code{x = 5} \\
        \addlinespace
        Kedekatan ke mesin & Langsung & Sangat dekat & Jauh \\
        \addlinespace
        Kendali perangkat keras & Penuh & Penuh & Terbatas \\
        \addlinespace
        Portabilitas & Rendah & Rendah & Tinggi \\
        \addlinespace
        Kemudahan penulisan & Sangat sulit & Sulit & Mudah \\
        \bottomrule
    \end{tabular}
\end{table}

\subsection{Mengapa Bahasa Assembly Masih Dipakai}

Bahasa tingkat tinggi lebih produktif untuk membangun aplikasi
umum. Meskipun demikian, bahasa assembly tetap memiliki tempat pada
pekerjaan yang menuntut kendali langsung atas mesin.

\begin{itemize}
    \item \textbf{Kendali perangkat keras.} Akses ke register, memori,
          dan port masukan/keluaran dapat dilakukan secara langsung
          tanpa lapisan abstraksi.

    \item \textbf{Ukuran dan kecepatan program.} Pada sistem dengan
          memori terbatas, program assembly dapat dibuat sangat kecil
          dan hemat instruksi.

    \item \textbf{Pemahaman cara kerja prosesor.} Assembly
          memperlihatkan bagaimana instruksi diambil, bagaimana data
          berpindah antara register dan memori, serta bagaimana
          keputusan diambil berdasarkan \eng{flag}.

    \item \textbf{Analisis program.} Pemahaman instruksi mesin
          membantu mahasiswa menelaah mengapa sebuah program berperilaku
          tertentu pada level perangkat keras.
\end{itemize}

\begin{infobox}
\vspace{1.1em}
Bahasa assembly bukan pengganti bahasa tingkat tinggi. Keduanya
merupakan alat dengan tujuan berbeda. Bahasa tingkat tinggi dipakai
ketika produktivitas dan pemeliharaan menjadi prioritas, sedangkan
assembly dipakai ketika kendali dan pemahaman level mesin diperlukan.
\end{infobox}

\subsection{Bahasa Assembly 8086 dalam Mata Kuliah Ini}

Mata kuliah ini memakai prosesor Intel 8086 dalam mode real 16-bit.
Perangkat praktik yang dipakai adalah TASM untuk meng-%
\eng{assemble} berkas \file{.ASM}, TLINK untuk menghasilkan berkas
\file{.EXE} atau \file{.COM}, DEBUG untuk memeriksa register dan
memori, serta emu8086 untuk menulis dan menjalankan program secara
terpandu. Alur kerja perangkat tersebut dibahas pada bab tentang
memulai assembly.

Pada bab ini, fokus pembelajaran bukan pada penulisan program utuh,
melainkan pada dua hal yang menjadi dasar seluruh bab berikutnya:
kedudukan bahasa assembly dan cara merepresentasikan bilangan.
Keduanya diperlukan sebelum mahasiswa menulis instruksi, menghitung
alamat memori, atau menafsirkan isi register.

\begin{aktivitasbox}
\vspace{1.1em}
Sebelum melanjutkan, lakukan dua aktivitas berikut.
\begin{enumerate}
    \item Tuliskan satu operasi sederhana, misalnya menambahkan dua
          bilangan, dalam bentuk kalimat bahasa Indonesia, lalu
          bayangkan bagaimana operasi itu dipecah menjadi langkah
          mesin.
    \item Sebutkan tiga pekerjaan yang menurut Anda lebih tepat
          dikerjakan dengan bahasa assembly daripada bahasa tingkat
          tinggi, beserta alasannya.
\end{enumerate}
\end{aktivitasbox}
